\section{8.1 Target System Vision}
The end goal is a closed-loop training system that accepts semantic commands and executes safe, accurate shots. Example command path:

\texttt{"right shoulder" -> select target point -> compute launch solution -> fire -> verify hit/miss -> adapt next shot}

This requires tight coupling between perception and actuation.

\section{8.2 Functional Blocks}
\subsection{Block A: Command Interface}
\begin{itemize}
  \item voice/text parser,
  \item intent classification (body part vs zone target),
  \item command validation and confirmation.
\end{itemize}

\subsection{Block B: Target Semantic Resolver}
\begin{itemize}
  \item map command token to 3D target in world frame,
  \item estimate confidence of selected body joint,
  \item gate on minimum confidence and camera support.
\end{itemize}

\subsection{Block C: Launcher Frame Transform}
\begin{itemize}
  \item transform target from world frame to launcher frame,
  \item apply launcher pose calibration,
  \item include static offsets (barrel origin, release point).
\end{itemize}

\subsection{Block D: Ballistics Solver}
Given target point \((x\_t, y\_t, z\_t)\) in launcher frame and release constraints, solve for initial velocity magnitude \(v\_0\) and azimuth/elevation angles \((\phi, \theta)\). For simplified projectile model (without spin and strong drag):

\[
x(t)=v_0\cos\theta\cos\phi \cdot t,
\]
\[
y(t)=v_0\cos\theta\sin\phi \cdot t,
\]
\[
z(t)=v_0\sin\theta\cdot t - \frac{1}{2}gt^2.
\]

In practice, launcher-specific calibration tables and empirical correction are needed due to wheel dynamics, spin, and drag.

\subsection{Block E: Actuation and Safety Controller}
\begin{itemize}
  \item send shot command,
  \item enforce no-fire constraints,
  \item require clear field and confidence threshold,
  \item emergency stop integration.
\end{itemize}

\subsection{Block F: Visual Verification and Adaptation}
\begin{itemize}
  \item reconstruct actual trajectory,
  \item compute impact/closest-approach to target,
  \item update correction model for future shots.
\end{itemize}

\section{8.3 Why Perception Quality Directly Affects Safety}
If 3D target estimate is wrong by 20-30 cm, a body-part-targeted shot can miss significantly and potentially create unsafe trajectories. Therefore, perception confidence gating is mandatory.

Recommended no-fire conditions:
\begin{itemize}
  \item fewer than 3 cameras support target joint,
  \item reprojection error above threshold,
  \item rapid target uncertainty spikes,
  \item second person detected in risk zone,
  \item stale calibration flag.
\end{itemize}

\section{8.4 Integration Phases}
\subsection{Table 8.1 Integration phases and exit criteria}
\begin{verbatim}
| Phase | Description | Exit criteria |
|---|---|---|
| P1 | Perception freeze and rigid GT qualification | rigid mean <= 60 mm, P95 <= 90 mm |
| P2 | Command parser and target semantic binding | >95% command parsing accuracy in test set |
| P3 | Ballistic solver in simulation | trajectory error < predefined tolerance |
| P4 | HIL supervised firing tests | zero safety violations, logged verification |
| P5 | Semi-autonomous training mode | stable performance across multi-session trials |
\end{verbatim}

\section{8.5 Voice Command Integration Notes}
Voice recognition can be implemented as modular front end:
\begin{itemize}
  \item offline keyword spotting or online ASR,
  \item strict vocabulary of target terms,
  \item explicit confirmation step before fire command.
\end{itemize}

Because command errors are high-risk, a two-step interaction is preferred:
\begin{enumerate}
  \item system repeats interpreted target,
  \item operator confirms.
\end{enumerate}

\section{8.6 Practical Engineering Constraints}
\begin{itemize}
  \item USB camera bandwidth and decode load can bottleneck live FPS.
  \item GPU acceleration is required for robust near-real-time operation.
  \item Camera mounts must be mechanically stable across sessions.
  \item Calibration updates must be part of standard operating procedure.
\end{itemize}

\section{8.7 Ethical and Safety Considerations}
The system should include:
\begin{itemize}
  \item operator override,
  \item activity logging for all commands and shots,
  \item explicit exclusion of bystanders,
  \item transparent uncertainty signaling,
  \item conservative default behavior under uncertainty.
\end{itemize}

A research prototype must never claim production safety without formal certification and controlled testing.

\section{8.8 Expected Impact}
If completed, this architecture can replace fixed-program training with adaptive, measurable drills. It can also reduce dependence on distributed physical sensors by using camera-defined target geometry and 3D event verification.


\section{8.9 Launcher Dynamics Calibration Requirements}
Even with perfect target coordinates, launcher output depends on wheel speed curves, spin-induced drift, release timing, and mechanical vibration. Therefore, ballistic planning should include empirical launcher calibration:
\begin{itemize}
  \item map command parameters to observed trajectory outcomes,
  \item fit correction surfaces for speed and angle,
  \item include uncertainty bounds in shot planning.
\end{itemize}

This calibration must be repeated after maintenance or hardware modifications.

\section{8.10 Control Strategy Options}
Three control strategies are feasible:

\begin{enumerate}
  \item Lookup-table control: fastest to deploy, uses pre-measured shot map.
  \item Model-based control: uses projectile equations with calibrated corrections.
  \item Hybrid adaptive control: starts with model/lookup and updates from visual feedback.
\end{enumerate}

For this project stage, the recommended path is model-based + empirical correction, then gradual adaptation.

\section{8.11 Human Factors and User Experience}
A practical training system must present clear operator feedback:
\begin{itemize}
  \item selected target and confidence score,
  \item whether system is in safe-to-fire state,
  \item reason for no-fire decisions,
  \item post-shot hit/miss and distance-to-target metrics.
\end{itemize}

Transparent feedback improves trust and enables coaches to interpret system behavior.

\section{8.12 Deployment Roadmap to Field Trials}
A realistic near-term roadmap:
\begin{enumerate}
  \item lock mounts and finalize rigid GT campaign,
  \item integrate command parser and target resolver,
  \item calibrate launcher-to-world transform,
  \item run low-energy supervised HIL shots,
  \item evaluate hit statistics and safety logs,
  \item expand to realistic training drills with controlled progression.
\end{enumerate}

This roadmap is feasible within phased research milestones and supports future publication-quality results.



\section{8.13 Interface Definition Between Perception and Launcher Controller}
A clean interface is recommended:

Input to controller:
\begin{itemize}
  \item \texttt{target\_point\_world\_mm},
  \item \texttt{target\_type} (joint/zone),
  \item \texttt{confidence\_score},
  \item \texttt{timestamp},
  \item \texttt{uncertainty\_score},
  \item \texttt{safety\_flags}.
\end{itemize}

Output from controller:
\begin{itemize}
  \item planned launch parameters,
  \item fire/no-fire decision,
  \item reason code,
  \item post-shot feedback summary.
\end{itemize}

This explicit contract prevents hidden assumptions between software modules.

\section{8.14 Progressive Safety Envelope}
A progressive safety envelope should be implemented:
\begin{enumerate}
  \item simulation-only mode,
  \item dry-run command mode (no firing),
  \item low-energy supervised firing,
  \item bounded autonomous mode.
\end{enumerate}

Each level should have quantitative pass/fail criteria.

\section{8.15 Validation Campaign for Launcher Stage}
A full launcher-stage validation campaign should include:
\begin{itemize}
  \item static target hit tests,
  \item moving target tracking without firing,
  \item supervised firing to large zones,
  \item supervised firing to body-part proxies,
  \item stress tests with occlusions and lighting shifts.
\end{itemize}

Each campaign must log confidence and no-fire decisions, not only successful hits.

\section{8.16 Research Publication Opportunities}
Potential publication directions after next stage:
\begin{itemize}
  \item practical robust extrinsics in constrained indoor arenas,
  \item protocol-driven benchmarking for sports CV systems,
  \item perception-aware safe actuation framework for training robotics.
\end{itemize}

The current thesis establishes the data and methods foundation for these outcomes.



\section{8.17 Data Logging Requirements for Autonomous Mode}
When launcher control is introduced, logs must include:
\begin{itemize}
  \item target command and parsed intent,
  \item selected 3D target and confidence,
  \item launch parameters sent,
  \item pre-fire safety checks,
  \item post-shot trajectory verification,
  \item hit/miss and residual distance.
\end{itemize}

Such logs are essential for safety audits and model improvement.

\section{8.18 Failure Recovery Strategy}
Controller should define explicit recovery policies:
\begin{itemize}
  \item if perception confidence drops before fire: cancel shot,
  \item if confidence drops after fire: mark low-confidence outcome,
  \item if repeated uncertainty events occur: enter safe hold mode,
  \item require operator acknowledgment to resume.
\end{itemize}

This state-machine approach prevents uncontrolled behavior.

\section{8.19 Calibration-Aware Command Scheduling}
Autonomous commands should be blocked when calibration age exceeds a threshold or when drift monitor flags inconsistency. Scheduling commands based on calibration health avoids accumulating hidden risk.

\section{8.20 Human-in-the-Loop Training Progression}
A recommended progression:
\begin{enumerate}
  \item coach-assisted command selection,
  \item system proposes target and waits confirmation,
  \item system executes low-energy shot,
  \item coach reviews visual verification,
  \item system updates adaptation parameters.
\end{enumerate}

This creates a safe transition from manual operation to increasing autonomy.

\section{8.21 Ethical Communication in Stakeholder Reporting}
It is important to communicate project status honestly:
\begin{itemize}
  \item perception system is validated with known limitations,
  \item launcher autonomy is planned and partially architected,
  \item safety-critical deployment requires additional testing.
\end{itemize}

This communication style aligns with research integrity and protects project credibility.

\section{8.22 Chapter Summary}
The architecture and control roadmap are technically grounded and immediately actionable for the next project phase. The design emphasizes measurable readiness and safety at every step.



\section{8.23 Proposed Software Stack for Controller Integration}
A practical software split is:
\begin{itemize}
  \item perception process (capture + inference + triangulation),
  \item decision process (target resolution + safety gating),
  \item control process (ballistic solve + actuator IO),
  \item monitoring process (visualization + logs + alerts).
\end{itemize}

Inter-process communication can use lightweight message queues with timestamped packets.

\section{8.24 Real-Time Validation Dashboard Requirements}
Before autonomous mode, operators need a dashboard displaying:
\begin{itemize}
  \item current target and confidence,
  \item camera participation per estimate,
  \item calibration health flag,
  \item no-fire reason codes,
  \item recent hit/miss history.
\end{itemize}

This dashboard is crucial for supervised transition stages.

\section{8.25 Human-in-the-Loop Ethics and Accountability}
In mixed human-machine training, accountability should be explicit. The system should preserve logs linking each fire command to perception state and operator approvals. This enables post-event analysis and responsible use in training contexts.

\section{8.26 Final Integration Perspective}
The integration task is not blocked by unknown fundamentals; it is blocked by engineering hardening steps already identified by this thesis. This is a strong position for next-phase execution.
